\documentclass[border=10pt]{standalone}
\usepackage[utf8]{inputenc}
\usepackage{nacal}
\usepackage{pgfplots}
\usepackage{tikz-3dplot}
\usepackage{tikz}
\usetikzlibrary{matrix,decorations.pathreplacing}
\usetikzlibrary{calc,angles,positioning,intersections,quotes,decorations.markings,arrows.meta}
\usepackage{tkz-euclide}

\begin{document}
  \def\Rojo{red!50!black}
  \def\RojoOscuro{red!20!black}
  \def\RojoClaro{red!90!black}
  \def\Verde{green!50!black}
  \def\VerdeOscuro{green!30!black}
  \def\VerdeClaro{green!50!gray}
  \def\Azul{blue!50!black}
  \def\AzulOscuro{blue!35!black}
  \def\AzulClaro{blue!60!gray}
  \def\Datos{green!40!gray}
  \def\Ajuste{blue!50!red}
  \def\Residuo{black!25!gray}
  \def\Regresor{blue!75!white}

  % theta = 35 grados, hipotenusa de longitud 4:
  %   cateto explicado = 4*cos(35) = 3.277 ; cateto residual = 4*sin(35) = 2.294
  \begin{tikzpicture}
    [scale=1.15,
     axis/.style={-,gray,very thin},
     vector/.style={-{Stealth[length=3mm, width=2mm]},very thick},
     guide/.style={dashed,thin,gray}]

    \coordinate (O)  at (0    , 0    );  % origen
    \coordinate (A)  at (3.277, 0    );  % ajuste en desviaciones
    \coordinate (B)  at (3.277, 2.294);  % datos en desviaciones
    \coordinate (Ex) at (1.400, 0    );  % el propio regresor en desviaciones
    \coordinate (Fin) at (3.55 , 0    );  % final del eje horizontal

    % la recta generada por el regresor en desviaciones; su etiqueta va al
    % extremo IZQUIERDO para dejar libre la derecha, donde ahora se rotula el
    % cateto explicado (antes ocupaba una fila entera bajo la figura)
    \draw[axis] (-0.55,0) -- (Fin);
    \node[anchor=north east,gray,yshift=-0.5mm] at (-0.55,0)
      {$\mathcal{L}(\boldsymbol{x}-\boldsymbol{\mathop{\overline{x}}})$};

    % el ángulo theta (el mismo de R^2=cos^2(theta) de la lección 8)
    \draw pic[-,draw=black,angle radius=26,\Verde,fill=\VerdeClaro!30] {angle = A--O--B};
    \node[\Verde!60!black,anchor=west,inner sep=1pt] at (17.5:0.95) {$\theta$};

    % ángulo recto entre el cateto explicado y el residuo
    \tkzMarkRightAngle[size=.28,gray,very thin,fill=lightgray!70,opacity=0.5](B,A,O)

    % el regresor en desviaciones, medido aparte bajo el eje: el cateto
    % explicado es un múltiplo suyo, y su longitud es la que divide a s
    \draw[guide] (O)  -- (0    ,-0.72);
    \draw[guide] (Ex) -- (1.400,-0.72);
    \draw[{Stealth[length=2mm,width=1.4mm]}-{Stealth[length=2mm,width=1.4mm]},\Regresor,thick]
      (0,-0.55) -- (1.400,-0.55);
    \node[\Regresor!70!black,anchor=west,xshift=1mm] at (1.400,-0.55)
      {$\boldsymbol{x}-\boldsymbol{\mathop{\overline{x}}}$};

    % los tres lados del triángulo
    \draw[vector,\Ajuste] (O) -- (A);
    \node[\Ajuste,anchor=north west,xshift=1.5mm,yshift=-0.5mm] at (A)
      {$\boldsymbol{\mathop{\widehat{y}}}-\boldsymbol{\mathop{\overline{y}}}
        =\hat\beta_2(\boldsymbol{x}-\boldsymbol{\mathop{\overline{x}}})$};
    \draw[vector,\Residuo] (A) -- (B)
      node[anchor=west,xshift=0.8mm]{$\boldsymbol{\mathop{\widehat{e}}}$};
    \draw[vector,\Datos] (O) -- (B)
      node[anchor=south east,xshift=-0.5mm,yshift=0.5mm]
      {$\boldsymbol{y}-\boldsymbol{\mathop{\overline{y}}}$};

    % las longitudes de los tres lados
    \node[\Ajuste,anchor=south,yshift=0.6mm] at (2.25,0) {$\sqrt{\mathrm{SEC}}$};
    \node[\Residuo!70!black,anchor=east,xshift=-0.8mm] at (3.277,1.147) {$\sqrt{\mathrm{SRC}}$};
    \node[\Datos!80!black,anchor=south east,xshift=-1mm,yshift=1mm] at (1.64,1.147) {$\sqrt{\mathrm{STC}}$};

  \end{tikzpicture}
\end{document}
